\section{5~V converter derivation}\label{sec:5v-converter}

The 5~V supply uses a TPS629203 converter (\cref{sec:tps629203}). This section
specifies its circuit requirements and external components.
\Cref{sec:bcd-driver} defines U1 to U4. \Cref{tab:5v-designators} defines the
remaining designators.

\begin{table}[htbp]
\centering
\caption{5~V converter: component designators}\label{tab:5v-designators}
\begin{tabular}{|l|l|}
\hline
\textbf{Designator} & \textbf{Function} \\
\hline
\des{L}{5V} & Buck inductor \\
\hline
\des{C}{5VIN} & Input capacitance \\
\hline
\des{C}{5VO1}, \des{C}{5VO2} & Output capacitance \\
\hline
\des{R}{5VMODE} & Device configuration resistor (MODE/S-CONF) \\
\hline
\des{R}{5VSET} & Output-voltage selection resistor (VSET) \\
\hline
C1 to C4 & U1 to U4 per-package bypass capacitors \\
\hline
\end{tabular}
\end{table}

\begin{figure}[htbp]
\centering
\includegraphics[width=\textwidth]{5v}
\caption{5~V converter: reference schematic}\label{fig:5v}
\end{figure}

\Cref{fig:5v} shows the connections; the four driver packages and their bypass
capacitors are not drawn.

The supply includes the converter, its input and output filters, and two
configuration resistors. It supplies the four K155ID1 $V_{CC}$ pins. C1 to C4
bypass these pins. The rail has no other load. The converter needs no bootstrap
capacitor (\cref{sec:tps629203}).

\subsection{Load budget}\label{sec:5v-load}

Each K155ID1 draws \SI{16}{\milli\ampere} typical and
\SI{25}{\milli\ampere} maximum (\cref{sec:bcd-driver}). Cathode current comes
from the 170~V rail. It returns through the selected open-collector output and
does not pass through $V_{CC}$. Thus, the 5~V load is only the driver logic
current:
\begin{align*}
  I_{\mathrm{load,typ}} &= 4 \times \SI{16}{\milli\ampere} = \SI{64}{\milli\ampere}, &
  I_{\mathrm{load,max}} &= 4 \times \SI{25}{\milli\ampere} = \SI{100}{\milli\ampere},
\end{align*}
\begin{align*}
  P_{\mathrm{out,typ}} &= 5 \times 0.064 = \SI{0.32}{\watt}, &
  P_{\mathrm{out,max}} &= 5 \times 0.100 = \SI{0.50}{\watt}.
\end{align*}
Bipolar TTL bias networks set the supply current. Output state has little effect
on it (inferred). Thus, the load is static for this analysis. C1 to C4 are
\SI{100}{\nano\farad} bypass capacitors. They add no DC load.

\subsection{Operating point}\label{sec:5v-input}

$V_{\mathrm{BUS}}$ connects directly to VIN and EN. The regulated input range is
\SIrange{8.55}{12.6}{\volt}, from nominal \SI{9}{\volt} and
\SI{12}{\volt} contracts. \Cref{sec:5v-limited} analyzes the
\SI{5}{\volt} contract separately. The regulated range is inside the
\SIrange{3}{17}{\volt} recommended range. The \SI{12.6}{\volt} maximum is
$70\,\%$ of the \SI{18}{\volt} absolute maximum. It meets the $80\,\%$ stress
limit used in this document.

The ideal buck duty cycle $D = V_{OUT}/V_{IN}$ spans
\begin{align*}
  D_{\min} &= \frac{5}{12.6} = 0.397 \quad (V_{IN}=\SI{12.6}{\volt}), &
  D_{\max} &= \frac{5}{8.55} = 0.585 \quad (V_{IN}=\SI{8.55}{\volt}).
\end{align*}
Conduction-loss corrections to $D$ are of order $I_{\mathrm{load}} R_{DS(\mathrm{on})}
\approx \SI{25}{\milli\volt}$ against a 5~V output (below 0.5\,\%) and are ignored
throughout. The converter reaches \SI{100}{\percent} duty only when
$V_{IN}\leq V_{OUT}+I_{OUT}(R_{DS(\mathrm{on})}+R_L)=\SI{5.04}{\volt}$.
The minimum regulated input is \SI{8.55}{\volt}.

\subsection{Output voltage and device configuration}\label{sec:5v-vset}

Two resistors program the converter. The converter reads both only during
start-up (\cref{sec:tps629203}).

\paragraph{Output voltage.} An \SI{80.6}{\kilo\ohm} resistor from VSET to ground
selects the internal \SI{5.0}{\volt} divider. Thus,
$\des{R}{5VSET} = \SI{80.6}{\kilo\ohm}$. Use an E96 $1\,\%$ resistor with a
temperature coefficient no more than $\pm200$~ppm/\si{\celsius}. The internal
divider keeps output tolerance at $\pm1\,\%$. An external divider would add
ratio tolerance. Connect VOS directly to the output-capacitor positive terminal.

\paragraph{Mode, frequency, and discharge.} MODE/S-CONF selects internal
feedback, automatic power save, efficiency enhancement, and output discharge.
The load is below the power-save threshold under all conditions
(\cref{sec:5v-op}). Output discharge gives a controlled ramp-down. Use
$\des{R}{5VMODE} = \SI{27.4}{\kilo\ohm}$. It has the same tolerance and
temperature-coefficient requirements as \des{R}{5VSET}.

\subsection{Inductor}\label{sec:5v-inductor}

The recommended inductance range is \SIrange{1.0}{4.7}{\micro\henry}. The
nominal value is \SI{2.2}{\micro\henry}. The datasheet does not recommend a
lower value when the inductor voltage is high. A lower value increases peak and
negative current. A higher value decreases ripple but can decrease efficiency
and transient response. The recommended filter table uses
\SI{2.2}{\micro\henry}. Thus, $\des{L}{5V} = \SI{2.2}{\micro\henry}$.

\paragraph{Ripple current.} In automatic power save with efficiency enhancement,
the on-time is set by the conversion ratio,
$t_{ON} = \SI{100}{\nano\second}\times V_{IN}/(V_{IN}-V_{OUT})$, and the ripple
follows from it:
\begin{align*}
  \Delta I_L &= \frac{V_{IN}-V_{OUT}}{L}\,t_{ON}
  &&\SI{8.55}{\volt}: \frac{3.55}{\SI{2.2}{\micro\henry}}\times\SI{241}{\nano\second}
     = \SI{0.39}{\ampere}, \\
  &&&\SI{12.6}{\volt}: \frac{7.6}{\SI{2.2}{\micro\henry}}\times\SI{166}{\nano\second}
     = \SI{0.57}{\ampere},
\end{align*}
These values agree with the datasheet curve. The ripple must not exceed
\SI{1.6}{\ampere}. The calculated maximum is \SI{0.57}{\ampere}, which has a
factor of 2.8 margin.

\paragraph{Saturation current.} The largest current the converter can drive into
the inductor is set by its own current limit, extended by the propagation delay
before the high-side switch actually opens:
\begin{align*}
  I_{\mathrm{peak}} &= I_{LIMH} + \frac{V_{IN}-V_{OUT}}{L}\times\SI{50}{\nano\second}
    = 1.4 + \frac{7.6}{\SI{1.76}{\micro\henry}}\times\SI{50}{\nano\second}
    = \SI{1.62}{\ampere},
\end{align*}
evaluated at the maximum high-side current limit and minimum inductance. Require
$I_{\mathrm{sat}}\geq\SI{1.7}{\ampere}$ and
${\rm DCR}\leq\SI{100}{\milli\ohm}$. The maximum operating peak is
\SI{0.57}{\ampere}. The current limit sets the saturation-current requirement.
At \SI{100}{\milli\ampere}, the DCR limit dissipates less than
\SI{1}{\milli\watt}.

\subsection{Conduction mode and switching frequency}\label{sec:5v-op}

With efficiency enhancement, switching frequency follows the conversion ratio:
\begin{align*}
  f_{SW} &= 10\,V_{OUT}\frac{V_{IN}-V_{OUT}}{V_{IN}^2}
  &&\SI{8.55}{\volt}: \SI{2.43}{\mega\hertz}; \quad
    \SI{10}{\volt}: \SI{2.50}{\mega\hertz}; \quad
    \SI{12.6}{\volt}: \SI{2.39}{\mega\hertz},
\end{align*}
The result spans \SIrange{2.39}{2.50}{\mega\hertz}. This analysis uses an
approximately \SI{2.4}{\mega\hertz} switching frequency.

The converter runs in power save mode whenever the load falls below half the
ripple current (\cref{sec:5v-inductor}):
\begin{align*}
  I_{\mathrm{load(PSM)}} &= \tfrac{1}{2}\Delta I_L
  &&\SI{8.55}{\volt}: \SI{194}{\milli\ampere}; \quad
    \SI{12.6}{\volt}: \SI{286}{\milli\ampere}.
\end{align*}
The \SIrange{64}{100}{\milli\ampere} load is below the power-save threshold at
all conditions. The device disables power save when input voltage is within
approximately $15\,\%$ of output. This threshold is \SI{5.75}{\volt}, below the
\SI{8.55}{\volt} minimum input. Therefore, the converter always uses power-save
operation in the regulated input range.

\subsection{Minimum on-time}\label{sec:5v-ontime}

The on-time in power save is the $t_{ON}$ of \cref{sec:5v-inductor},
\SI{241}{\nano\second} at the minimum input falling to \SI{166}{\nano\second} at
the maximum, against a \SI{40}{\nano\second} floor: a margin of \num{4.1} at the
binding corner. In forced pulse-width modulation the on-time would be
$(V_{OUT}/V_{IN})/f_{SW} = \SI{166}{\nano\second}$ at the same corner, the same
figure to within the frequency's own variation. The complementary
constraint, a minimum off-time of about \SI{80}{\nano\second}, binds at the
opposite end: at \SI{8.55}{\volt} input and \SI{2.43}{\mega\hertz} the off-time is
$1/f_{SW} - t_{ON} = \SI{171}{\nano\second}$, a margin of \num{2.1}. Neither
constrains the design, and the off-time is the tighter of the two only because
the duty cycle is highest at minimum input.

\subsection{Current-limit margin}\label{sec:5v-ilim}

The minimum high-side current limit is \SI{0.8}{\ampere}. The maximum commanded
power-save peak is \SI{0.57}{\ampere}:
\begin{align*}
  \frac{I_{LIM}^{\min}}{I_{\mathrm{peak}}} &= \frac{0.8}{0.57} = 1.4.
\end{align*}
Propagation delay increases the effective minimum trip current to
\SI{0.97}{\ampere}. This gives a factor of 1.7 margin. If the current reaches
the limit, the high-side switch turns off. The low-side switch then decreases
the inductor current. This protection operates in each cycle and does not latch.

\subsection{Output capacitance and rail accuracy}\label{sec:5v-cout}

The part requires effective output capacitance between \SIrange{10}{100}{\micro\farad}
(\cref{sec:tps629203}), and the recommended filter table pairs the
\SI{2.2}{\micro\henry} inductance with nominal values from \SIrange{10}{100}{\micro\farad}.
The analysis assumes that X5R or X7R capacitors retain at least $50\,\%$ of
nominal capacitance at \SI{5}{\volt} (inferred). Use
$\des{C}{5VO1} = \des{C}{5VO2} = \SI{22}{\micro\farad}$. The pair gives
\SI{44}{\micro\farad} nominal and \SI{22}{\micro\farad} effective. One
\SI{22}{\micro\farad} part would give only \SI{11}{\micro\farad} effective.
This value has little margin above the \SI{10}{\micro\farad} minimum.

\paragraph{Ripple.} Output ripple in power save mode is
\begin{align*}
  \Delta V &= \frac{L\,V_{IN}^2}{200\,C}
    \left(\frac{1}{V_{IN}-V_{OUT}}+\frac{1}{V_{OUT}}\right)
  &&\SI{8.55}{\volt}: \SI{17.6}{\milli\volt}; \quad
    \SI{12.6}{\volt}: \SI{26.3}{\milli\volt},
\end{align*}
with the inductance in \si{\micro\henry} and the effective capacitance in
\si{\micro\farad}. Ripple grows with input voltage, so the maximum bus voltage is
the binding corner.

\paragraph{Accuracy stack against the load's window.} The internal divider holds
the output to $\pm1\,\%$ over the full junction-temperature range, so regulation
alone spans \SIrange{4.95}{5.05}{\volt}. Adding the worst-case ripple to both
ends:
\begin{align*}
  V_{\mathrm{rail}} &\in [\,4.95-0.027,\; 5.05+0.027\,]
    = [\,\SI{4.92}{\volt},\ \SI{5.08}{\volt}\,]
\end{align*}
against the K155ID1's recommended \SIrange{4.75}{5.25}{\volt} supply window
(\cref{sec:bcd-driver}): approximately \SI{170}{\milli\volt} of margin on each
side. This is the figure the internal divider protects. An external divider built
from 1\,\% resistors would add roughly $\pm1.7\,\%$ of ratio error on top of the
reference's own tolerance, consuming about half the margin above for no
benefit, since \SI{5.0}{\volt} is already one of the preset options.

\subsection{Input capacitance}\label{sec:5v-cin}

The device requires at least \SI{3}{\micro\farad} effective input capacitance.
At \SI{12.6}{\volt}, a \SI{25}{\volt} ceramic can lose approximately half its
capacitance (inferred). Therefore, use
$\des{C}{5VIN} = \SI{10}{\micro\farad}$. It gives approximately
\SI{5}{\micro\farad} effective. The maximum input RMS ripple current occurs
near a duty cycle of 0.5:
\begin{align*}
  I_{\mathrm{RMS}} &\approx I_{\mathrm{out}}\sqrt{D(1-D)}
    \approx 0.100\sqrt{0.5 \times 0.5} = \SI{50}{\milli\ampere},
\end{align*}
The calculated ripple is \SI{50}{\milli\ampere}. A ceramic
\SI{10}{\micro\farad} capacitor can typically carry this current (inferred).
Use a \SI{25}{\volt} rating to reduce DC-bias derating.

\subsection{Enable, soft start, and dissipation}\label{sec:5v-ss}

EN connects directly to $V_{\mathrm{BUS}}$. Its \SI{18}{\volt} rating matches
VIN. No divider is necessary. The maximum rising UVLO threshold is
\SI{3}{\volt}, below the \SI{8.55}{\volt} minimum input. The power-good output is
unused and unconnected.

Start-up is internally timed: up to \SI{1.8}{\milli\second} of delay from enable
to the first switching pulse in the internal-divider configuration, then a
soft-start ramp of up to \SI{700}{\micro\second} to the target. The charging
current that ramp implies is
\begin{align*}
  I_{\mathrm{chg}} &= C_{\mathrm{out}}\frac{dV}{dt}
    = \SI{22}{\micro\farad}\times\frac{\SI{5}{\volt}}{\SI{600}{\micro\second}}
    = \SI{183}{\milli\ampere},
\end{align*}
which even summed with the full \SI{100}{\milli\ampere} load stays a factor of
\num{2.8} below the minimum current limit: no inrush interaction with the limit.
The 5~V rail reaches regulation approximately \SI{2.5}{\milli\second} after
$V_{\mathrm{BUS}}$. This converter adds no cross-rail sequence requirement.

Dissipation, using the conservative $\eta = 0.85$ carried for the low-voltage
rails:
\begin{align*}
  P_{\mathrm{diss}} &= P_{\mathrm{out,max}}\left(\frac{1}{\eta}-1\right)
    = 0.50 \times 0.176 = \SI{88}{\milli\watt}, &
  \Delta T &= \theta_{JA} P_{\mathrm{diss}} = 120 \times 0.088 \approx \SI{11}{\kelvin},
\end{align*}
so $T_J \approx \SI{51}{\celsius}$ at a \SI{40}{\celsius} ambient, far below the
\SI{125}{\celsius} recommended junction limit and further still below the
\SI{170}{\celsius} thermal shutdown. The \SI{120}{\celsius\per\watt} figure is
the datasheet's JEDEC-board value; its evaluation-board figure is
\SI{60}{\celsius\per\watt}, so the result above is the pessimistic of the two.
No thermal management beyond the footprint's own copper is required.

\subsection{Operation on a \SI{5}{\volt} contract}\label{sec:5v-limited}

A \SI{5}{\volt} contract gives $V_{\mathrm{BUS}}$ from
\SIrange{4.75}{5.5}{\volt}. This range is below the regulated input range. The
converter enters 100-percent-duty operation. The high-side switch stays on, and
the output follows the input minus
\begin{align*}
  \Delta V &= I_{\mathrm{load}}\left(R_{DS(\mathrm{on})}+R_L\right)
    = 0.1\,(0.25+0.10) = \SI{35}{\milli\volt}
\end{align*}
at the \SI{100}{\milli\ampere} bound. The rail therefore lands between
\SI{4.72}{\volt} and \SI{5.47}{\volt} and is unregulated, tracking whatever the
bus does.

In limited operation, the high-voltage supply stays off. The drivers then have
no cathode current. In full operation, the drivers switch cathode current from
the tracking 5~V rail. At minimum input, the rail is approximately
\SI{30}{\milli\volt} below the recommended \SI{4.75}{\volt} limit. This small
undervoltage changes TTL levels but does not cause electrical stress (inferred).
At maximum input, the rail reaches \SI{5.47}{\volt}. This is $78\,\%$ of the
\SI{7}{\volt} absolute maximum. Thus, only the rail-accuracy result in
\cref{sec:5v-cout} requires a \SI{9}{\volt} or \SI{12}{\volt} contract.

\subsection{Summary}\label{sec:5v-summary}

The analysis uses two inferred values. Ceramic capacitors retain at least
\SI{50}{\percent} of nominal capacitance at their operating bias. Converter
efficiency is 0.85. The selected inductor must have
$I_{\mathrm{sat}}\geq\SI{1.7}{\ampere}$ and
${\rm DCR}\leq\SI{100}{\milli\ohm}$.

{\footnotesize
\begin{longtable}{|p{0.23\textwidth}|p{0.335\textwidth}|p{0.335\textwidth}|}
\caption{5~V converter: specification summary}\label{tab:5v-converter-summary}\\
\hline
\textbf{Parameter} & \textbf{Derived requirement} & \textbf{Specification / result} \\
\hline
\endfirsthead
\caption[]{5~V converter: specification summary (continued)}\\
\hline
\textbf{Parameter} & \textbf{Derived requirement} & \textbf{Specification / result} \\
\hline
\endhead
\hline
\endlastfoot
Load current & Not applicable & $4\times I_{CC}$: \SI{64}{\milli\ampere} typ /
  \SI{100}{\milli\ampere} max, static \\
\hline
Input window & inside \SIrange{3}{17}{\volt} recommended; \SI{80}{\percent} of
  the \SI{18}{\volt} absolute maximum & regulation across
  \SIrange{8.55}{12.6}{\volt}; tracking in dropout across
  \SIrange{4.75}{5.5}{\volt} (\cref{sec:5v-limited}); \SI{70}{\percent} of
  rating at \SI{12.6}{\volt} \\
\hline
Duty cycle & Not applicable & 0.397 to 0.585; \SI{100}{\percent} mode unreachable above
  \SI{5.04}{\volt} input \\
\hline
Output voltage \des{R}{5VSET} & \SI{5.0}{\volt} preset; E96 1\,\%,
  $\pm200$~ppm/\si{\celsius} & \SI{80.6}{\kilo\ohm} on VSET \\
\hline
Configuration \des{R}{5VMODE} & internal divider, auto power save with efficiency
  enhancement, output discharge & \SI{27.4}{\kilo\ohm} on MODE/S-CONF \\
\hline
Inductance \des{L}{5V} & \SIrange{1.0}{4.7}{\micro\henry} recommended; not below
  \SI{2.2}{\micro\henry} at high input & \SI{2.2}{\micro\henry}; ripple
  \SIrange{0.39}{0.57}{\ampere}, inside the \SI{1.6}{\ampere} caution \\
\hline
Inductor $I_{\mathrm{sat}}$, DCR & $\geq\SI{1.62}{\ampere}$ from the
  delay-extended current limit; DCR $\leq\SI{100}{\milli\ohm}$ &
  $\geq\SI{1.7}{\ampere}$; part to be selected \\
\hline
Switching frequency & Not applicable & \SIrange{2.39}{2.50}{\mega\hertz} across the input
  window \\
\hline
Conduction mode & Not applicable & load below the \SIrange{194}{286}{\milli\ampere}
  power-save boundary at every corner \\
\hline
Minimum on-time & $\geq\SI{40}{\nano\second}$ on; $\geq\SI{80}{\nano\second}$
  off & \SI{166}{\nano\second} on ($4.1\times$); \SI{171}{\nano\second} off
  ($2.1\times$) \\
\hline
Current-limit margin & Not applicable & \SI{0.57}{\ampere} peak vs \SI{0.8}{\ampere}
  minimum limit; \num{1.7} including propagation delay \\
\hline
Output capacitance \des{C}{5VO1}, \des{C}{5VO2} & \SIrange{10}{100}{\micro\farad} effective &
  $2\times\SI{22}{\micro\farad}$: \SI{22}{\micro\farad} effective \\
\hline
Output ripple & Not applicable & \SI{26.3}{\milli\volt} at the maximum input \\
\hline
Rail accuracy at load & \SIrange{4.75}{5.25}{\volt} (K155ID1) &
  \SIrange{4.92}{5.08}{\volt} including ripple, on a \SI{9}{\volt} or
  \SI{12}{\volt} contract; \cref{sec:5v-limited} otherwise \\
\hline
Input capacitance \des{C}{5VIN} & $\geq\SI{3}{\micro\farad}$ effective;
  $\geq\SI{16}{\volt}$; RMS $\geq\SI{49}{\milli\ampere}$ &
  \SI{10}{\micro\farad}; \SI{25}{\volt} preferred \\
\hline
Soft start & Not applicable & internal; $\leq\SI{2.5}{\milli\second}$ delay plus ramp;
  inrush \SI{183}{\milli\ampere} \\
\hline
EN configuration & terminated, $\leq\SI{18}{\volt}$ & tied to
  $V_{\mathrm{BUS}}$ \\
\hline
Converter dissipation & Not applicable & \SI{88}{\milli\watt}, $\Delta T \approx
  \SI{11}{\kelvin}$ (SOT-5X3, $\theta_{JA}=120$~\si{\celsius\per\watt}) \\
\hline
Bypass C1 to C4 & \SI{100}{\nano\farad} per TTL package & \SI{100}{\nano\farad}
  $\times 4$ \\
\hline
\end{longtable}
}

The inductor saturation rating remains open because no part is selected. Select
a part that meets the stated value. The converter current limit sets the
inductor requirement. The specified inductor and capacitors also comply with
the manufacturer filter table.

\subsection{Component selection}\label{sec:5v-parts}

\Cref{tab:5v-parts} gives the component-selection criteria. Body size is a
requirement only for the fixed converter package. Other body sizes are typical
values that meet the electrical requirements.

{\footnotesize
\begin{longtable}{|L{0.10\textwidth}|L{0.14\textwidth}|L{0.11\textwidth}|L{0.33\textwidth}|L{0.12\textwidth}|}
\caption{5~V converter: component selection}\label{tab:5v-parts}\\
\hline
\textbf{Designator} & \textbf{Class} & \textbf{Value} & \textbf{Required ratings and characteristics} & \textbf{Body} \\
\hline
\endfirsthead
\caption[]{5~V converter: component selection (continued)}\\
\hline
\textbf{Designator} & \textbf{Class} & \textbf{Value} & \textbf{Required ratings and characteristics} & \textbf{Body} \\
\hline
\endhead
\hline
\endlastfoot
Not applicable & Synchronous buck converter & TPS629203 & VSET configuration selected, giving $\pm1\,\%$ output regulation over the junction range without an external divider & SOT-5X3 (SOT-583), $1.6\times2.1$~\si{\milli\meter} \\
\hline
\des{L}{5V} & Power inductor & \SI{2.2}{\micro\henry} & $I_{\mathrm{sat}}\geq\SI{1.7}{\ampere}$; DCR $\leq\SI{100}{\milli\ohm}$; $\pm20\,\%$ tolerance acceptable, the extremes staying inside the part's \SIrange{1.0}{4.7}{\micro\henry} recommended range; core rated for \SI{2.5}{\mega\hertz} & $\leq5\times5$~\si{\milli\meter}, $\leq\SI{4.5}{\milli\meter}$ high \\
\hline
\des{C}{5VIN} & Ceramic, input & \SI{10}{\micro\farad} & $\geq\SI{16}{\volt}$, \SI{25}{\volt} preferred; X5R or X7R; $\geq50\,\%$ capacitance retention at \SI{12.6}{\volt} bias; ripple current $\geq\SI{50}{\milli\ampere}$ RMS & 0805 \\
\hline
\des{C}{5VO1} & Ceramic, output & \SI{22}{\micro\farad} & $\geq\SI{10}{\volt}$; X5R or X7R; $\geq50\,\%$ capacitance retention at \SI{5}{\volt} bias & 1206 \\
\hline
\des{C}{5VO2} & Ceramic, output & \SI{22}{\micro\farad} & as \des{C}{5VO1} & 1206 \\
\hline
\des{R}{5VMODE} & Film resistor & \SI{27.4}{\kilo\ohm} & E96 1\,\%; temperature coefficient no worse than $\pm200$~ppm/\si{\celsius}; selects internal feedback, automatic power save with efficiency enhancement, and output discharge & 0402 \\
\hline
\des{R}{5VSET} & Film resistor & \SI{80.6}{\kilo\ohm} & E96 1\,\%; temperature coefficient no worse than $\pm200$~ppm/\si{\celsius}; selects the \SI{5.0}{\volt} preset & 0402 \\
\hline
C1 to C4 & Ceramic, TTL bypass & \SI{100}{\nano\farad} & $\geq\SI{10}{\volt}$ under the \SI{80}{\percent} derating rule; X7R; one per K155ID1 package & 0402 \\
\hline
\end{longtable}
}
